\section{Evaluation is easier than generation：可扩展监督的杠杆}

前面的方法依赖人类比较。要把它扩展到更强模型，必须回答一个关键问题：人类能否评价自己不会完成的任务？讲者的出发点是一个常见但非绝对的非对称性---生成高质量结果可能很难，看到候选后识别哪一个更好往往容易一些。

\subsection{从命题到直觉}

课程先用单独标题页停顿，是因为这句话承担后半讲的全部逻辑。如果 evaluation 与 generation 同样困难，那么人类一旦被模型能力超越就失去监督入口；若评价更容易，人类便可以借助候选、证据和批评，把监督能力延伸到自己无法独立完成的任务。本节先把这条原则写成有条件的成本优势，再通过四个类比检查它何时成立、何时会因隐藏证据或 evaluator 被操纵而失效。

\lecturefigure{slide-11-evaluation-title.jpg}{后半讲的核心原则：evaluation is easier than generation。}{Stanford Online 官方视频 00:20:19--00:20:32。}

\begin{importantbox}{评价优势的条件化表述}
更准确的说法是：给定候选结果、任务规范和足够可验证证据时，比较或局部检查的成本可能低于从零生成。优势来自搜索空间被候选压缩，而不是 evaluator 自动拥有真理。若错误隐藏在不可见状态、证据不可得、评价者可被说服或任务目标本身含糊，评价优势会减弱甚至消失。
\end{importantbox}

\subsection{四个类比应该怎样读}

Slide 列出 P/NP、观看比赛、购买产品和学术审稿。它们共享一个结构：执行者需要在巨大行动空间中找到方案，评价者看到少量候选后只需比较结果或指出局部问题。但这些例子不是证明 ``所有真实任务都可高效验证''；它们只是解释为什么 preference comparison 可能比 demonstration 更可扩展。

\lecturefigure{slide-12-evaluation-easier-examples.jpg}{评价比生成容易的类比：复杂性、体育、消费选择与论文评审。}{Stanford Online 官方视频 00:20:32--00:22:25。}

\begin{knowledgebox}{读图：把类比压成一个共同结构}
P/NP 强调找到解与验证证书的差异；体育强调观众能判断胜负却无法达到职业水平；消费选择强调比较体验不要求制造产品；审稿强调指出缺陷不等于写出更好的论文。共同点是 evaluator 获得了候选或证据。不同点是现实评价常有主观目标、隐藏变量和不完整证据，不能直接套用复杂性理论结论。
\end{knowledgebox}

若生成成本记为 $C_{\mathrm{gen}}(T)$，在辅助信息 $a$ 下的评价成本记为 $C_{\mathrm{eval}}(T\mid a)$，scalable oversight 希望构造 $a$ 使

\[
C_{\mathrm{eval}}(T\mid a) \ll C_{\mathrm{gen}}(T),
\]

其中，$T$ 表示任务，$a$ 可以是候选答案、引用、测试、critique、对话或工具结果。符号 $\ll$ 是设计目标，不是对所有任务成立的定理。

\begin{warningbox}{评价者也会被优化}
一旦模型知道什么会获得高分，它就不只优化任务，还会优化 evaluator 的观察过程。流畅表达、选择性证据、迎合偏见甚至隐藏缺陷都可能提高评分。评价更容易只有在评价流程本身对这些策略足够鲁棒时才转化为可靠训练信号。
\end{warningbox}

\subsection{本章小结}

评价优势为 scalable oversight 提供入口：人类不必独立生成专家级结果，也可能比较候选或检查具体证据。但这个优势依赖可见证据、明确规范与抗操纵的评价流程；它不是 ``人类永远能监督更强 AI'' 的保证。

